\documentclass[11pt]{article}
\usepackage[margin=1in]{geometry}
\usepackage{amsmath, amssymb}
\usepackage{enumitem}
\usepackage{tcolorbox}
\usepackage{xcolor}

\title{\textbf{Question 4b: The Simple Version} \\ \large You've Got This!}
\author{}
\date{}

\begin{document}

\maketitle

\section{Don't Panic - Here's What's Really Being Asked}

Question 4b sounds complicated, but it's actually asking something very intuitive:

\begin{tcolorbox}[colback=blue!10!white,colframe=blue!75!black]
\textbf{The Real Question:}

``I know there were $n$ arrivals by time $s$. What can I say about arrivals at some OTHER time $t$?''

That's it. That's the whole question.
\end{tcolorbox}

The trick is: your answer depends on whether $t$ is BEFORE or AFTER time $s$.

\section{Let's Use a Concrete Story}

Imagine customers arriving at a coffee shop. On average, $\lambda = 3$ customers per hour arrive.

\vspace{0.3cm}

\textbf{The Fact You Know:} By 9am (time $s = 9$), exactly 13 customers have arrived ($N_9 = 13$).

\vspace{0.3cm}

Now the question asks: \textit{``What about arrivals at time $t$?''}

\section{Case 1: Looking BACKWARD (When $t < s$)}

\subsection{The Question}

You're at 9am. You know 13 customers have arrived by now.

\textbf{Question:} How many of those 13 customers had arrived by 4am?

\begin{verbatim}
Time:    midnight    4am         9am
         |-----------|-----------|
                     ?           13 customers total
\end{verbatim}

\subsection{The Answer: Binomial}

Think about it: You have 13 customers. Each one arrived at some random time between midnight and 9am.

\textbf{For each customer:}
\begin{itemize}
    \item 4am is $\frac{4}{9}$ of the way from midnight to 9am
    \item So each customer has a $\frac{4}{9}$ chance of arriving before 4am
    \item Like flipping a coin with probability $\frac{4}{9}$ for each of the 13 customers
\end{itemize}

\begin{tcolorbox}[colback=green!10!white,colframe=green!75!black,title=Answer for Case 1]
\[
(N_4 \mid N_9 = 13) \sim \text{Binomial}\left(13, \frac{4}{9}\right)
\]

\textbf{Translation:} Out of 13 customers, each has a $\frac{4}{9}$ chance of being in the first 4 hours.

\textbf{Mean:} On average, $13 \times \frac{4}{9} \approx 5.78$ customers arrived by 4am.

\textbf{Notice:} We DON'T need to know $\lambda = 3$! The rate doesn't matter when looking backward.
\end{tcolorbox}

\subsection{Why This Makes Sense}

You already know the 13 customers arrived somewhere between midnight and 9am. You're just asking: ``How are they split between [midnight-4am] and [4am-9am]?''

Since arrivals are random, they split proportionally: $\frac{4}{9}$ vs $\frac{5}{9}$.

\section{Case 2: Same Time (When $t = s$)}

\subsection{The Question}

You're at 9am. You know 13 customers have arrived.

\textbf{Question:} How many customers arrived by 9am?

\subsection{The Answer}

Obviously 13. You just told me 13 customers arrived by 9am!

\begin{tcolorbox}[colback=blue!10!white,colframe=blue!75!black,title=Answer for Case 2]
\[
N_9 = 13 \quad \text{(for certain)}
\]

This case is trivial - if $t = s$, then $N_t = n$ with 100\% certainty.
\end{tcolorbox}

\section{Case 3: Looking FORWARD (When $t > s$)}

\subsection{The Question}

You're at 9am. You know 13 customers have arrived so far.

\textbf{Question:} How many customers will have arrived by 2pm (time $t = 14$)?

\begin{verbatim}
Time:    midnight    9am         2pm
         |-----------|-----------|
         13 customers    ???
\end{verbatim}

\subsection{The Answer: n + Poisson}

Think about it:
\begin{itemize}
    \item You DEFINITELY have the 13 customers who already arrived
    \item PLUS however many MORE arrive between 9am and 2pm (5 hours)
    \item Those future arrivals follow a Poisson with rate $\lambda \times 5$ hours
\end{itemize}

\begin{tcolorbox}[colback=purple!10!white,colframe=purple!75!black,title=Answer for Case 3]
\[
(N_{14} \mid N_9 = 13) = 13 + \text{Poisson}(\lambda \times 5) = 13 + \text{Poisson}(3 \times 5) = 13 + \text{Poisson}(15)
\]

\textbf{Translation:} 13 known customers PLUS a Poisson(15) random number of future customers.

\textbf{Mean:} On average, $13 + 15 = 28$ customers by 2pm.

\textbf{Notice:} NOW we need $\lambda = 3$! We need to know the rate to predict the future.
\end{tcolorbox}

\subsection{Why This Makes Sense}

You can't change the past - you know 13 customers already came. But the future is random - you expect about 15 more customers (since $\lambda = 3$ per hour $\times$ 5 hours = 15), but it could be more or less.

\section{Side-by-Side Comparison}

\begin{table}[h]
\centering
\begin{tabular}{|l|c|c|}
\hline
\textbf{Aspect} & \textbf{Case 1: $t < s$ (Past)} & \textbf{Case 3: $t > s$ (Future)} \\
\hline
Example & 4am vs 9am & 9am vs 2pm \\
\hline
Known info & 13 by 9am & 13 by 9am \\
\hline
Question & How many by 4am? & How many by 2pm? \\
\hline
Distribution & Binomial$(13, \frac{4}{9})$ & $13 +$ Poisson$(15)$ \\
\hline
Need $\lambda$? & NO & YES \\
\hline
Mean & $13 \times \frac{4}{9} \approx 5.78$ & $13 + 15 = 28$ \\
\hline
Intuition & Split 13 proportionally & Known + future \\
\hline
\end{tabular}
\end{table}

\section{The General Formula (What 4b Wants)}

For ANY times $t$ and $s$:

\begin{tcolorbox}[colback=yellow!10!white,colframe=orange!75!black,title=Complete Answer]
Given $N_s = n$, the distribution of $N_t$ is:

\vspace{0.3cm}

\textbf{If $t < s$ (looking at the past):}
\[
N_t \sim \text{Binomial}\left(n, \frac{t}{s}\right)
\]

\vspace{0.3cm}

\textbf{If $t = s$ (same time):}
\[
N_t = n \quad \text{(certain)}
\]

\vspace{0.3cm}

\textbf{If $t > s$ (looking at the future):}
\[
N_t = n + \text{Poisson}(\lambda(t - s))
\]
\end{tcolorbox}

\section{How to Actually Solve Problems}

\subsection{Step 1: Identify Which Case}

Look at $t$ and $s$. Is $t < s$, $t = s$, or $t > s$?

\subsection{Step 2: Use the Right Formula}

\begin{itemize}
    \item Past ($t < s$): Binomial - use $\frac{t}{s}$ as the probability
    \item Same ($t = s$): Just $n$
    \item Future ($t > s$): $n$ + Poisson - need $\lambda$
\end{itemize}

\subsection{Step 3: Calculate}

Plug in numbers and compute!

\section{Practice Problem}

\textbf{Given:} $\lambda = 2$ customers per hour. By time $s = 5$ hours, $n = 8$ customers arrived.

\subsection{Question A: How many by time $t = 2$ hours?}

\textbf{Identify:} $t = 2 < s = 5$ (PAST)

\textbf{Formula:} Binomial$(8, \frac{2}{5})$

\textbf{Mean:} $8 \times \frac{2}{5} = 3.2$ customers

\textbf{Full answer:} $(N_2 \mid N_5 = 8) \sim \text{Binomial}(8, 0.4)$

\subsection{Question B: How many by time $t = 10$ hours?}

\textbf{Identify:} $t = 10 > s = 5$ (FUTURE)

\textbf{Formula:} $8 + \text{Poisson}(\lambda(t-s)) = 8 + \text{Poisson}(2 \times 5) = 8 + \text{Poisson}(10)$

\textbf{Mean:} $8 + 10 = 18$ customers

\textbf{Full answer:} $(N_{10} \mid N_5 = 8) = 8 + \text{Poisson}(10)$

\section{Common Mistakes (and How to Avoid Them)}

\subsection{Mistake 1: Using Binomial for Future}

\textcolor{red}{\textbf{Wrong:}} ``By 9am there were 13 customers. By 2pm: Binomial$(13, \frac{14}{9})$''

\textbf{Why wrong:} You can't have probability $> 1$! Also, you can't have fewer than 13 customers at 2pm.

\textcolor{green!60!black}{\textbf{Right:}} Use $13 + \text{Poisson}$ for future times.

\subsection{Mistake 2: Forgetting $\lambda$ for Future}

\textcolor{red}{\textbf{Wrong:}} ``I don't need $\lambda$ for any case''

\textbf{Why wrong:} For the future, you MUST know the rate to predict new arrivals.

\textcolor{green!60!black}{\textbf{Right:}} Past doesn't need $\lambda$, but future does!

\subsection{Mistake 3: Using Wrong Time Interval}

\textcolor{red}{\textbf{Wrong:}} Future case: $n + \text{Poisson}(\lambda t)$

\textbf{Why wrong:} Should be Poisson$(\lambda(t - s))$ - the time from $s$ to $t$, not from 0 to $t$.

\textcolor{green!60!black}{\textbf{Right:}} Use the DIFFERENCE: $t - s$

\section{You've Got This!}

\begin{tcolorbox}[colback=green!20!white,colframe=green!75!black]
\textbf{The Big Picture:}

\begin{itemize}
    \item This is just asking: ``Given info at time $s$, what about time $t$?''
    \item Past = split what you know (Binomial)
    \item Future = known + random new arrivals (Poisson)
    \item It's actually intuitive once you think about the coffee shop story!
\end{itemize}

\vspace{0.3cm}

\textbf{How to succeed:}
\begin{enumerate}
    \item Always identify: is $t$ before or after $s$?
    \item Use the right formula
    \item For past: just need $n$ and $\frac{t}{s}$
    \item For future: need $n$, $\lambda$, and $(t-s)$
\end{enumerate}

\vspace{0.3cm}

\textbf{You can do this.} The concept is actually simpler than it looks. Just remember the coffee shop!
\end{tcolorbox}

\end{document}
